\begin{textsample}{POS Dim 2 – df – Score 27.00 – df/df\_em\_57.txt}  \label{ex:f2_pos_152}
ORGANIZAÇÃO, FORMAS E MODALIDADES DE \textbf{OFERTA} A EPT articulada ao Ensino Médio se dá por meio de diferentes estratégias de formação nas \textbf{Instituições} Educacionais da rede de ensino ou em \textbf{Instituições} Parceiras devidamente autorizadas, com os \textbf{cursos} FIC ou de \textbf{Qualificação} Profissional e \textbf{Cursos} \textbf{Técnicos} de Nível Médio, com organização curricular própria, de acordo com as \textbf{legislações} \textbf{vigentes}.
Ademais, segundo Cordão ( 2010 ) : > E a educação profissional vai além dos \textbf{cursos} \textbf{técnicos} de nível médio integrados ou articulados com o ensino médio.
Ela integra todo o \textbf{itinerário} formativo das pessoas, preparando -as para o mundo do trabalho, para definirem seus próprios \textbf{itinerários} de profissionalização ( CORDÃO, 2010, p. 37 ). \textbf{QUALIFICAÇÃO} PROFISSIONAL OU FORMAÇÃO INICIAL E CONTINUADA ( FIC ) Os \textbf{cursos} de \textbf{Qualificação} Profissional ou FIC são organizados e \textbf{ofertados} conforme denominação, \textbf{carga} \textbf{horária}, \textbf{escolaridade} e perfis profissionais definidos segundo o \textbf{Guia} Nacional de \textbf{Cursos} FIC, editado pelo MEC, e/ou que constem da Classificação Brasileira de Ocupações ( CBO ), editada pelo Ministério do Trabalho ( BRASIL, 2002b ), além de se constituírem como partes integrantes da matriz curricular dos \textbf{cursos} \textbf{técnicos} de nível médio, conforme possibilidades de saídas intermediárias, \textbf{previstas} no CNCT.
Após a conclusão do \textbf{curso}, o estudante recebe certificado de acordo com as normativas do DF e a devida descrição da denominação do \textbf{Guia} Nacional de \textbf{Cursos} FIC e/ou da CBO.
A elaboração do Plano de \textbf{Curso} para \textbf{oferta} desses \textbf{cursos} deve ocorrer consoante as normativas definidas pela área responsável no âmbito do Sistema de Ensino do Distrito Federal, sendo necessariamente construída de forma articulada com a Formação Geral Básica. \textbf{CURSOS} \textbf{TÉCNICOS} DE NÍVEL MÉDIO Os \textbf{cursos} \textbf{técnicos} são \textbf{destinados} a estudantes matriculados no Ensino Médio ou egressos, conforme preconiza o \textbf{Catálogo} Nacional de \textbf{Cursos} \textbf{Técnicos} ( CNCT ) do Ministério da Educação.
O \textbf{Catálogo} Nacional de \textbf{Cursos} \textbf{Técnicos} é instrumento que disciplina a \textbf{oferta} de \textbf{cursos} de educação profissional \textbf{técnica} de nível médio, para orientar as \textbf{instituições}, os estudantes e a sociedade em geral.
É um referencial para subsidiar o planejamento dos \textbf{cursos} e correspondentes \textbf{qualificações} profissionais e \textbf{especializações} de nível médio ( BRASIL, 2014b, p. 8 ).
O CNCT, em sua última atualização, disponibiliza a lista de \textbf{cursos} \textbf{técnicos}, os quais estão agrupados em eixos tecnológicos, sendo que cada perfil profissional exige uma \textbf{carga} \textbf{horária} mínima.
O acesso e a permanência nos \textbf{cursos} promovem a preparação \textbf{técnica} para a inserção dos estudantes no mundo do trabalho, sem, no entanto, deixar de possibilitar a continuidade dos estudos para a \textbf{Especialização} \textbf{Técnica} de Nível Médio e para o Ensino Superior.
Durante a formação, os estudantes têm aulas teóricas e práticas, as quais podem ser realizadas em laboratórios ou em empresas, a fim de aproximar a formação das experiências com o perfil profissional.
A elaboração do Plano de \textbf{Curso} para a \textbf{oferta} de \textbf{cursos} \textbf{técnicos} de nível médio deve obedecer às normativas e \textbf{legislações} \textbf{vigentes}, sendo necessariamente construída de forma articulada com a Formação Geral Básica.
FORMAS DE \textbf{OFERTA} - Integrada ao Ensino Médio : a formação ocorre em um mesmo \textbf{currículo} que integra a Base Nacional Comum Curricular e as unidades curriculares específicas do \textbf{curso} \textbf{técnico} \textbf{ofertado}, com matrícula única do estudante e, portanto, na mesma \textbf{Instituição} Educacional. - Concomitante ao Ensino Médio : a formação ocorre com \textbf{currículos} distintos e articulados na forma, integrando a Base Nacional Comum Curricular que é \textbf{cursada} em uma \textbf{Instituição} Educacional e as unidades curriculares específicas do \textbf{curso} \textbf{técnico}, \textbf{ofertadas} na mesma \textbf{Instituição} Educacional ou em outra \textbf{Instituição}, sendo para isso necessário o estabelecimento de duas matrículas.
Em ambas as formas de \textbf{oferta}, o \textbf{currículo} deverá ser executado de forma integrada, permitindo que o estudante \textbf{curse} mais de um Itinerário Formativo, simultânea ou sequencialmente.
POSSIBILIDADES DE \textbf{OFERTA} - Forma Presencial : a organização curricular \textbf{prevê} a execução de todas as unidades curriculares e seus respectivos conteúdos por meio de aulas em que os estudantes e os docentes estão fisicamente, na maioria das vezes, no mesmo local e ao mesmo tempo, o que exige o cumprimento da \textbf{carga} \textbf{horária} \textbf{prevista} com definição de hora de início e término. - Modalidade de Educação a Distância ( EaD ) : a mediação didático-pedagógica nos processos de ensino e aprendizagem ocorre com a utilização de meios e tecnologias de informação e comunicação, nos quais estudantes e docentes, na maioria das vezes, desenvolvem atividades educativas em lugares ou tempos diversos.
REGISTRO EM SISTEMA NACIONAL DE INFORMAÇÕES DA EPT A fim de garantir a validade nacional dos diplomas expedidos e registrados na própria \textbf{instituição} educacional, todo \textbf{curso} \textbf{técnico} de nível médio, com Plano de \textbf{Curso} aprovado pelo CEDF, deve ser registrado no Sistema Nacional de \textbf{Cursos} \textbf{Técnicos} \textbf{vigente}, definido pelo MEC.
FORMALIZAÇÃO DE PARCERIAS PARA \textbf{OFERTA} DO \textbf{ITINERÁRIO} DE EPT A \textbf{oferta} dos \textbf{Itinerários} Formativos no Novo Ensino Médio possibilita a articulação da SEEDF com \textbf{instituições} parceiras, na \textbf{garantia} e efetivação de direitos, e estimula o desenvolvimento de uma \textbf{gestão} pública democrática e participativa.
Nesse sentido, as parcerias se apresentam como uma forma de aproximar as \textbf{políticas} públicas dos cidadãos e das realidades locais, possibilitando a solução conjunta de problemas no cenário educacional.
Em relação à \textbf{certificação}, caberá às \textbf{instituições} parceiras emitir certificados de \textbf{qualificação} profissional ou outros documentos comprobatórios das atividades concluídas sob sua responsabilidade, sendo que os documentos comprobatórios deverão ser incorporados pela \textbf{instituição} de origem no Histórico Escolar do estudante.
Firmar parcerias, para a implementação do Novo Ensino Médio, objetiva reunir agentes no cumprimento de sua função social, para a consecução de finalidades de interesse público, a fim de \textbf{atender} às necessidades e expectativas dos estudantes, fortalecer seu interesse, engajamento e protagonismo, bem como promover a sua permanência e aprendizagem.

% matched lemmas: atender, carga, catálogo, certificação, currículo, cursar, curso, destinar, escolaridade, especialização, garantia, gestão, guia, horário, instituição, itinerário, legislação, oferta, ofertar, política, prever, qualificação, técnico, vigente
\end{textsample}
